\documentclass[11pt,a4paper]{article}
\usepackage[margin=2.5cm]{geometry}
\usepackage[T1]{fontenc}
\usepackage{lmodern}
\usepackage{enumitem}
\usepackage{titlesec}
\titleformat{\section}{\large\bfseries}{\thesection.}{0.5em}{}

\title{Narrative Photo Selection and Sequencing}
\author{\textbf{Project Proposal}}
\date{}

\begin{document}
\maketitle

\section{Problem Statement}

Many users have large collections of unordered digital photographs.
Existing tools summarize photos with timestamps or text captions.
But real photo albums do not always follow chronological order.

This project develops an automated system that selects $N$ photos from a
collection of $K$ photos, orders them and aims for maximum coverage and coherent order.
The system must satisfy two goals:

\begin{enumerate}
  \item Select photos that have high visual quality and broad coverage of the wider dataset.
  \item Arrange the selected photos in a coherent order that makes sense to a human viewer.
\end{enumerate}

E.g.\ a photo of a beach followed by a photo of a snowy mountain into another photo of a beach would be a bad sequence.

\section{Implementation Plan}

\begin{enumerate}
  \item Explore, Find and Collect appropriate datasets for the problem.
  \begin{itemize}
    \item The dataset should be large enough to have a good evaluation.
    \item The dataset should have a wide range of photos.
    \item The dataset should have photos with high visual quality as well as ones with low visual quality.
    \item It should be separated by subjects -- e.g.\ user1 went on a trip to Japan and has 150 photos, user2 has 200 photos from a birthday party\ldots
    \item Ideally it should be labeled in some way -- this is not a requirement.
  \end{itemize}

  \item Develop a system that can select $N$ photos from the dataset.
  \begin{itemize}
    \item The system should be able to select photos that are not in chronological order.
    \item The system should be able to select photos that are in chronological order.
    \item The system should be ``tunable'' in multiple ways so that selection can be matched to user prefences.
  \end{itemize}

  \item Evaluate the system.
  \begin{itemize}
    \item Evaluate the system on the dataset. (if labeled)
    \item If needed survey users to evaluate the system.
  \end{itemize}

  \item Write a report outlining the results and discoveries.
\end{enumerate}

\end{document}
